QtRocket is an open-source model-rocket flight simulator: a C++ simulation core
shared by a Qt6 desktop application, a command-line REPL called
\code{qtrocket-cli}, and a standalone 3D design viewer called
\code{qtrocket-visualizer}. This guide teaches the command-line interface ---
today the most feature-complete way to drive QtRocket, and the only way to
\emph{author} rocket designs. Everything you build here saves to \code{.qrd}
files that the GUI and the visualizer open directly.

The guide is organized as a short concepts primer
(\cref{sec:getting-started,sec:concepts}), five hands-on tutorials
(\cref{sec:tut1,sec:tut2,sec:tut3,sec:tut4,sec:tut5}), and a complete command
reference (\cref{sec:reference}). The tutorials are the heart of it: each one
builds something real, flies it, and reads the results, introducing commands as
they earn their place.

\begin{note}[Every transcript in this guide is real]
The session transcripts in this guide are not typed-up approximations: each one
is the captured output of the released \code{qtrocket-cli} binary running the
tutorial scripts included alongside this document (\cref{app:repro}). The
screenshots are captures of \code{qtrocket-visualizer} displaying the
\code{.qrd} files those same sessions saved, and the flight plots are drawn
from the CSV files they wrote. Where a long payload is shortened, the elision
is marked in \textcolor{qrGray}{\textit{[bracketed gray italics]}} --- nothing
else is edited. If you run the scripts, you will get these numbers.
\end{note}

\paragraph{Reading the transcripts.} Lines you type are shown in
\textcolor{qrBlue}{\textbf{bold blue}}. The CLI answers every command with a
line starting \ok{} (success) or \err{} (failure); those and the occasional
\warnline{WARN} are colored the way you should read them. Muted
\muted{\#~gray lines} are comments and banners the CLI prints or ignores.

\paragraph{Conventions.} All quantities are SI --- meters, kilograms, seconds,
Newtons; densities are kg/m$^3$. Altitude is the $z$~axis with the ground at
$z=0$. Motor designations follow the standard hobby impulse classes (A, B, C,
\dots), where each letter doubles the total impulse of the one before it.
